\section{Prompt templates}
\label{app:prompts}

All TTMG LLM calls route through the same chat interface (OpenAI-compatible).
The three module-specific prompts below are what the writer, linker, and
retriever query-parser actually see; small differences between the
session-batched writer and the single-turn writer are flagged.

\subsection{Session-batched writer}
\label{app:writer_prompt}
\begin{quote}\small
\emph{System:} You extract factual claims from conversation turns for a
long-term memory system. Only extract information the speaker ACTUALLY
asserts. Do not invent content. Respond with strict JSON only.\\[2pt]
\emph{User:} Extract factual claims from the conversation below.
Session metadata: \texttt{session\_id}, \texttt{session\_ts}.
Turns are numbered. For each turn containing factual content (facts,
preferences, plans, states), output claims with fields \texttt{turn\_id},
\texttt{speaker}, \texttt{content}, \texttt{subject}, \texttt{predicate},
\texttt{object}, \texttt{valid\_from}, \texttt{valid\_to},
\texttt{polarity} $\in\{\texttt{assert},\texttt{deny}\}$,
\texttt{volatility} $\in\{\texttt{stable},\texttt{preference},
\texttt{state}\}$, \texttt{confidence} $\in [0,1]$. Skip greetings,
small talk and pure questions asked of the model. Keep content under 25
words. Return \texttt{\{"claims":[\ldots]\}}.
\end{quote}

\subsection{Link labeller}
\label{app:linker_prompt}
\begin{quote}\small
\emph{System:} You label pairs of factual claims as one of
\texttt{support}, \texttt{contradict}, \texttt{supersede}, or
\texttt{unrelated}. Respond in strict JSON only.\\[2pt]
\emph{User:} For the NEW claim and each CANDIDATE claim output a label
and a confidence in $[0,1]$.
\textbf{support}: same fact; both coexist. \textbf{contradict}:
disagree on the same $(s,p)$ at overlapping times. \textbf{supersede}:
the NEW claim replaces the CANDIDATE (same $(s,p)$, later
\texttt{valid\_from}, or explicit correction). \textbf{unrelated}:
otherwise. Be conservative: default to \texttt{unrelated} unless
subject and predicate match. Return
\texttt{\{"pairs":[\{"id":\ldots,"label":\ldots,"confidence":\ldots\}]\}}.
\end{quote}

\subsection{Query parser}
\label{app:parser_prompt}
\begin{quote}\small
\emph{System:} You parse a question asked of a long-term memory system.
Respond with strict JSON only.\\[2pt]
\emph{User:} Parse the question. Return fields
\texttt{temporal\_anchor} (time string or \texttt{null}),
\texttt{asks\_history} (question asks about a past state that may have
been superseded), \texttt{asks\_current} (question asks about the
current state), \texttt{entity} (main subject or \texttt{null}),
\texttt{intent} $\in\{\texttt{lookup}, \texttt{update\_check},
\texttt{temporal\_reasoning}, \texttt{small\_talk}\}$.
\end{quote}

\section{Hyper-parameters and inference-time router}
\label{app:hparams}

\begin{table}[h]
  \centering
  \caption{TTMG default hyper-parameters, used unchanged across all
  benchmarks.}
  \label{tab:hyperparams}
  \small
  \begin{tabular}{ll}
    \toprule
    Component & Value / model \\
    \midrule
    Embedding model & \texttt{all-MiniLM-L6-v2} \\
    Reader (answer generation) & \texttt{deepseek-v3.2} \\
    Writer / linker / parser & \texttt{Kimi-K2} \\
    Intent classifier (TTMG+R-LLM) & \texttt{Kimi-K2}, zero-shot \\
    Retrieval $k$ & 8 \\
    $k_{\text{keep}}$ at read time & 3 \\
    Linker candidate $k$ at write time & 3 \\
    Similarity gate for linker & 0.65 \\
    Bypass threshold & 0.85 \\
    Hard-edge threshold $\tau_{\text{hard}}$ & 0.7 \\
    Minimum claim-query similarity & 0.35 \\
    Session-batched writer & on \\
    Max sessions ingested per question & 3 \\
    Evaluation seed & 0 (stratified) \\
    \bottomrule
  \end{tabular}
\end{table}

\subsection{LLM intent classifier for TTMG+R-LLM}
\label{app:intent}
The TTMG+R-LLM router uses a one-shot Kimi-K2 call per incoming
question to predict one of the six LongMemEval types. The classifier
prompt lists the six labels with one example each and asks for a
single-key JSON response. On the 150-question stratified evaluation
slice, its overall accuracy is \IntentAccOverall\%. Per-class
recall is: \texttt{single-session-user} \IntentRecallSSU\%,
\texttt{single-session-assistant} \IntentRecallSSA\%,
\texttt{temporal-reasoning} \IntentRecallTR\%,
\texttt{multi-session} \IntentRecallMS\%, \texttt{knowledge-update}
39.1\%, \texttt{single-session-preference} 22.2\%. The high recall
on SSU and SSA is what makes the simple
``route-SSU-to-TTMG-else-Flat'' rule practical; the lower recall on
MS and KU is the dominant source of the gap to the oracle upper
bound and a natural target for a better-trained classifier.

\subsection{Reproducibility statement}
\label{app:reproduce}
We release the full source (writer/linker/parser/reader prompts,
the retriever algorithm, the claim schema, and all experiment
scripts), the N=150 stratified subsample identifiers and seed
(\texttt{seed=0}, stratified on \texttt{question\_type}), all
hyper-parameters (Table~\ref{tab:hyperparams}), and the raw
per-question outputs
(\texttt{results/pilot\_n150\_*.json}), which includes the
\texttt{question\_id}, \texttt{correct} flag, and metric counters
per question. McNemar $p$-values and bootstrap CIs are computed by
\texttt{scripts/analyze\_n150.py}. The LLM intent classifier is
implemented in \texttt{scripts/llm\_intent\_classifier.py}. All
calls route through the MAAS OpenAI-compatible endpoint; model
identifiers (\texttt{deepseek-v3.2}, \texttt{Kimi-K2}) are frozen
strings provided by the hosting platform.

\section{Dataset details}
\label{app:datasets}

LongMemEval-S contains 500 questions across six types:
\emph{single-session-user} (70), \emph{single-session-assistant}~(56),
\emph{single-session-preference}~(30), \emph{multi-session}~(133),
\emph{knowledge-update}~(78), and \emph{temporal-reasoning}~(133). A
further 30 questions bearing an \texttt{\_abs} suffix in their
question identifier are \emph{abstention} probes: the reference answer
states that the information was never mentioned, so the correct
response is ``I do not know''. Each question is paired with an
average of $\sim\!53$ haystack sessions drawn from a user's
conversation history; \texttt{answer\_session\_ids} labels the
session(s) from which the answer may be derived.

Each of the 50 questions retrieves memory from its own haystack; we
cap the ingested haystack at the 3 highest-priority sessions (the
answer session plus the two earliest matching sessions) to keep the
per-question ingestion budget manageable while preserving the ratio
of relevant-to-irrelevant sessions the LongMemEval benchmark intends.
